\documentclass{article}

\usepackage[english]{babel}
\usepackage[a4paper,top=2cm,bottom=2cm,left=3cm,right=3cm,marginparwidth=1.75cm]{geometry}
\usepackage{float}

% Useful packages
\usepackage{amsmath}
\usepackage{graphicx}
\usepackage[colorlinks=true, allcolors=blue]{hyperref}


\title{Kalea Corven}
\author{Karcag Tamas}

\begin{document}
\maketitle

\begin{abstract}
    Your abstract.
\end{abstract}

\section{Appearance}

\begin{figure}[H]
    \centering
    \includegraphics[width=0.6\textwidth]{src/Kalea.png}
    \caption{\label{fig:kalea}Kalea portrait.}
\end{figure}

Kalea Corven is a twenty-one-year-old woman of slim build, with striking blue eyes and unusual grey hair. Her appearance has always made her stand out, something that became a source of discomfort for her family from an early age.

She wears dark travelling clothes reinforced with practical leather armour, accompanied by sturdy leather boots. A long, worn black cloak with a damaged hood covers most of her clothing and allows her to conceal her appearance when necessary. Its condition reflects the hardships of life on the road rather than the comfortable existence she once knew.

Her clothes are chosen for practicality, protection, and anonymity. She avoids luxurious accessories and anything that might reveal her noble origins. Nevertheless, certain gestures, habits, and mannerisms occasionally betray the education and upbringing she received.

\section{Behaviour}

Kalea is intelligent, curious, and analytical. She has a strong desire to learn and continuously seeks opportunities to expand her knowledge and improve her abilities. She prefers to observe and understand a situation before taking action, often considering several possible outcomes before committing to a decision.

Her years spent in the family library made her particularly fond of books, research, and intellectual challenges. She is comfortable with written knowledge and theoretical discussions, although her limited experience of ordinary life sometimes leaves her less prepared for practical situations than she would like to admit.

Her noble upbringing has made her familiar with the customs and behaviour of the upper classes. However, she no longer identifies with the values associated with that world. She strongly dislikes arrogance, exploitation, and the abuse of authority, particularly when those in positions of power take advantage of people who cannot defend themselves.

Although she has ambitious ideals, Kalea understands that changing society cannot happen overnight. She is willing to pursue smaller, achievable goals if they bring tangible improvements to people's lives.

She is cautious about revealing personal information and does not easily trust strangers. Her current circumstances require discretion, and she knows that a careless word or an unexpected encounter could expose her identity.

Despite her careful nature, she can sometimes be stubborn when confronted with injustice. Her convictions may lead her to take risks even when remaining uninvolved would be the safer choice.

\section{Background}

\subsection{Noble Origins and Family}

Kalea was born the youngest child of House Corven, a noble family with an established position in high society. She
has an older brother and an older sister, both of whom were raised to fulfil the expectations and responsibilities associated with their family's standing.

Unlike her siblings, Kalea was never truly welcomed into the family. Her unusual grey hair, present from birth, made her an object of discomfort and embarrassment. Her parents considered her appearance undesirable and feared that it might attract unwanted attention or damage the family's reputation.

As a result, Kalea grew up largely isolated from the outside world. While her brother and sister were educated, introduced to influential people, and prepared for their future roles, she was kept away from public gatherings and family events. Her presence was rarely acknowledged outside the household.

Her family provided her with little affection or attention. Instead, she was mostly confined to the family library. Officially, this was presented as an opportunity to pursue her education in peace. In reality, it was also a convenient way to keep her out of sight.

The library became both her refuge and her prison.

She spent much of her childhood and adolescence surrounded by books, studying history, politics, philosophy, and the structure of society. With few opportunities to interact with people, books became her primary source of knowledge about the world beyond the estate.

Over time, she developed a sharp intellect and a particular interest in understanding how societies function, how power is distributed, and why some people are granted privileges while others are denied even basic opportunities.

Her family may have intended to keep her hidden, but in doing so, they gave her the means to question the very system that protected them.

\subsection{The Turning Point}

As Kalea grew older, she began to understand that her isolation was not temporary. Her family had no intention of allowing her to take part in public life or to enjoy the same opportunities as her siblings. She was expected to remain out of sight, dependent on the household, and grateful for the limited comforts she had been given.

The decisive moment came when she overheard a private conversation between members of her family. They were discussing her future, not as a daughter or an individual with her own wishes, but as an inconvenience that needed to be dealt with. Her continued presence was considered a liability, and any decision concerning her life would be made according to what best served the family's interests.

The conversation made something painfully clear: her family did not intend to give her a life of her own. They intended to keep controlling her, just as they always had.

Around the same time, her studies led her to discover records and correspondence that revealed more about the family's dealings with those beneath their social position. She began to see how the privileges enjoyed by the nobility were maintained through the exploitation and suffering of ordinary people.

The contrast was impossible to ignore. Her family had hidden her because they considered her appearance unacceptable, while outside their walls countless people were suffering because of a social structure her family actively benefited from.

For the first time, Kalea understood that her dissatisfaction was about more than her own treatment. She no longer wanted simply to escape her family's control. She wanted to challenge the world that allowed such treatment to be considered normal.

She decided that she could not remain there.

\subsection{The Escape}

Kalea did not make her decision impulsively. She knew that leaving a noble household would be difficult, and that her family had the resources and connections to search for her if they discovered what she had done.

She therefore spent time preparing in secret.

The library, which had confined her for most of her life, also provided her with the knowledge needed to plan her escape. She studied the estate's layout, learned the household's daily routines, and paid attention to the movements of servants, guards, and visitors. She waited until she understood when the estate was least attentive and which routes offered the best chance of leaving without being noticed.

She gradually gathered a small amount of money and a few essential belongings. She chose practical clothing over the fine garments prepared for someone of her birth, and obtained a dark travelling cloak to conceal her appearance. She took only what she could carry without drawing attention.

On the night she finally left, she waited until the household had settled. Using a service passage she had learned about from her observations, she avoided the main corridors and the rooms where members of the family and their guests were likely to be found.

She made her way through the less frequently used parts of the estate and eventually reached the grounds beyond the main residence. From there, she continued towards an outer exit used by servants and deliveries.

Every unfamiliar sound made her hesitate. She had spent her entire life inside those walls, and although she had studied the outside world extensively, she had never experienced it for herself.

Nevertheless, she kept moving.

By the time her absence was discovered, Kalea was already gone. She left behind the comfort and security of her birthright, along with the identity her family had assigned to her.

She knew that her escape would not remain a secret forever. Her family might eventually discover where she had gone, and their influence could make hiding difficult. From that moment onwards, remaining anonymous became essential to her survival.

\subsection{A Different Purpose}

Life outside the estate proved far more difficult than anything Kalea had learned from books. She had knowledge and education, but little practical experience of living among ordinary people. She had to learn how to find work, manage limited resources, recognise danger, and judge the intentions of strangers.

She also began to encounter people whose lives were very different from the privileged world she had known. For the first time, she could see the consequences of the inequalities she had previously studied from a distance.

These experiences strengthened her convictions.

Kalea's long-term ambition is to change the rigid social structure that divides people according to birth, wealth, and status. She believes that a person's worth should not be determined by the family they were born into or the position they occupy in society.

She wants to help create a world where people are treated as equals, where privilege does not grant someone the right to exploit others, and where those without power have the opportunity to improve their lives.

For now, however, such an ambition is far beyond her immediate reach. She therefore focuses on smaller, achievable objectives: helping oppressed people, supporting those without power or influence, and holding members of the privileged classes accountable when they abuse their position.

She continues to study and train, determined to turn her education into something useful. Knowledge alone cannot change the world, but she believes it can be the beginning of meaningful change.

Kalea has escaped the life that was chosen for her. Now she must discover what she is willing to risk to build the life, and the society, she believes in.

\end{document}